\documentclass[11pt,a4paper]{article}
\usepackage[italian]{babel}
\usepackage[margin=2.5cm]{geometry}
\usepackage{parskip}
\usepackage{amsmath, amssymb, amsthm}
\usepackage{booktabs}
\usepackage{array}
\usepackage{xcolor}
\usepackage{needspace}
\usepackage{enumitem}
\usepackage{listings}
\usepackage{tikz}
\usepackage[most]{tcolorbox}
\usepackage[hidelinks]{hyperref}
\usetikzlibrary{decorations.pathreplacing, shapes.geometric, arrows.meta, positioning}

% Niente vedove/orfane: i paragrafi non lasciano righe isolate a fine/inizio pagina
\widowpenalty=10000
\clubpenalty=10000

% Elenchi più compatti
\setlist{itemsep=2pt, parsep=0pt, topsep=4pt, partopsep=0pt}

% --- Riquadri con bordo nero, senza numero (si spezzano tra due pagine) ---
% Uso: \begin{definizione} ... \end{definizione}
%      \begin{teorema}[Nome] ... \end{teorema}  (il nome è facoltativo)
\tcbset{nota/.style={enhanced jigsaw, breakable, colback=white,
  colframe=black, boxrule=0.5pt, sharp corners}}
\NewTColorBox{definizione}{}{nota,
  before upper={\textbf{Definizione.}\ }}
\NewTColorBox{teorema}{o}{nota,
  before upper={\textbf{Teorema\IfValueT{#1}{ (#1)}.}\ }}
\NewTColorBox{proposizione}{o}{nota,
  before upper={\textbf{Proposizione\IfValueT{#1}{ (#1)}.}\ }}
\NewTColorBox{proprieta}{o}{nota,
  before upper={\textbf{Proprietà\IfValueT{#1}{ (#1)}.}\ }}
\NewTColorBox{assioma}{o}{nota, boxrule=1pt,
  before upper={\textbf{Assioma\IfValueT{#1}{ (#1)}.}\ }}

% --- Ambienti semplici (senza riquadro) ---
% Il titolo sta in cima al blocco, anche se il contenuto inizia con un riquadro o
% con codice; e non resta mai da solo in fondo a una pagina.
% Uso: \begin{esempio}[Titolo facoltativo] ... \end{esempio}
\newcommand{\newrunin}[2]{%
  \NewDocumentEnvironment{#1}{o}{%
    \par\Needspace{4\baselineskip}\addvspace{\medskipamount}\noindent
    \textbf{#2}\IfValueT{##1}{ (##1)}\textbf{.}\ \ignorespaces}%
  {\par\addvspace{\medskipamount}}}
\newrunin{esempio}{Esempio}
\newrunin{esercizio}{Esercizio}
\newrunin{osservazione}{Osservazione}

% V e F nelle tavole di verità
\newcommand{\V}{\textsf{V}}
\newcommand{\F}{\textsf{F}}

% --- Codice C ---
% Uso: \begin{codice} ... \end{codice}      codice C numerato
%      \begin{comando} ... \end{comando}    comandi da terminale
%      \begin{risultato} ... \end{risultato} output di un programma
%      \lstinline|printf("ciao");|           codice dentro al testo
\lstdefinestyle{c}{
  language=C,
  basicstyle=\ttfamily\small,
  keywordstyle=\color{blue}\bfseries,
  commentstyle=\color{gray}\itshape,
  stringstyle=\color{red!70!black},
  numbers=left,
  numberstyle=\tiny\color{gray},
  numbersep=8pt,
  columns=flexible,
  keepspaces=true,
  breaklines=true,
  showstringspaces=false,
  tabsize=4,
  morekeywords={include, define},
  % lettere accentate nei commenti
  literate={à}{{\`a}}1 {è}{{\`e}}1 {é}{{\'e}}1 {ì}{{\`i}}1
           {ò}{{\`o}}1 {ù}{{\`u}}1 {È}{{\`E}}1 {À}{{\`A}}1
}
\lstset{style=c}
\newtcblisting{codice}{listing only, nota, left=6mm,
  listing options={style=c}}
\newtcblisting{comando}{listing only, nota,
  listing options={basicstyle=\ttfamily\small, numbers=none, language={},
    columns=flexible, keepspaces=true, breaklines=true}}
\newtcblisting{risultato}{listing only, enhanced jigsaw, breakable,
  colback=black!5, colframe=black, boxrule=0.5pt, sharp corners,
  listing options={basicstyle=\ttfamily\small, numbers=none, language={},
    columns=flexible, keepspaces=true, breaklines=true}}

% --- Diagrammi di flusso (simboli standard) ---
\tikzset{
  terminale/.style = {ellipse, draw, minimum width=2.5cm, minimum height=0.9cm},
  io/.style        = {trapezium, trapezium left angle=70, trapezium right angle=110,
                      draw, minimum width=2.5cm, minimum height=0.9cm},
  processo/.style  = {rectangle, draw, minimum width=2.5cm, minimum height=0.9cm},
  decisione/.style = {diamond, draw, aspect=2, minimum width=2.5cm},
  freccia/.style   = {thick, -{Stealth}}
}

\title{Esempi sugli estremi e insiemi non limitati}
\author{Marco Cerbella}
\date{Lezione 3 -- 1 ottobre 2026}

\begin{document}
\maketitle

\section{Un altro esempio: $B = \left\{ \dfrac{n}{n+1} \mid n \in \mathbb{N} \right\}$}

\begin{esempio}
Sia $B = \{x_n = \frac{n}{n+1},\ n \in \mathbb{N}\} = \{0, \frac12, \frac23, \frac34, \dots\}$. Poiché $x_n \leq 1$ per ogni $n$ (infatti $n \leq n + 1$) e $x_n \geq 0$, si ha $B \subseteq [0, 1]$: $B$ è limitato. In particolare, essendo limitato superiormente, per (R4) esiste $\sup B$.

\textbf{Minimo.} $x_0 = 0 \in B$ e $x_n \geq 0$ per ogni $x_n \in B$: quindi $0 = \min B$.

\textbf{Estremo superiore.} Dimostriamo che $\sup B = 1$, cioè che:
\begin{enumerate}
    \item $x_n \leq 1$ per ogni $x_n \in B$ (già visto);
    \item per ogni $\varepsilon > 0$ esiste $x_n \in B$ tale che $1 - \varepsilon < x_n$.
\end{enumerate}
Per assurdo neghiamo la 2: esiste $\bar\varepsilon > 0$ tale che $1 - \bar\varepsilon \geq x_n = \frac{n}{n+1}$ per ogni $n \in \mathbb{N}$. Allora
\[
(n+1)(1 - \bar\varepsilon) \geq n \iff n + 1 - \bar\varepsilon n - \bar\varepsilon \geq n
\iff 1 - \bar\varepsilon \geq \bar\varepsilon n
\iff n \leq \frac{1 - \bar\varepsilon}{\bar\varepsilon} \quad \forall\, n \in \mathbb{N}.
\]
Questo è assurdo, perché $n$ può essere grande a piacere. Quindi $1 = \sup B$.

Poiché $1 \notin B$ (si ha $x_n < 1$ per ogni $n$), $B$ non ha massimo.
\end{esempio}

\section{Insiemi non limitati}

Un insieme $A$ non sempre ammette maggioranti o minoranti.

\begin{definizione}
Si pone
\begin{itemize}
    \item $\sup A = +\infty$ se $A$ non è limitato superiormente, cioè
          $\forall\, L \ \exists\, x \in A : x > L$;
    \item $\inf A = -\infty$ se $A$ non è limitato inferiormente, cioè
          $\forall\, \ell \ \exists\, x \in A : x < \ell$.
\end{itemize}
\end{definizione}

\begin{esempio}
\leavevmode
\begin{itemize}
    \item $A = [0, +\infty)$ è limitato inferiormente ma non superiormente;
    \item $A = (-\infty, 0]$ è limitato superiormente ma non inferiormente;
    \item $[-2, +\infty)$ è limitato inferiormente ma non superiormente;
    \item $(-\infty, 5]$ è limitato superiormente ma non inferiormente;
    \item $\,]2, 1000[$ è limitato: $2 = \inf\,]2, 1000[$ e
          $1000 = \sup\,]2, 1000[$, ma l'insieme non ha né minimo né massimo
          (perché $2, 1000 \notin \,]2, 1000[$);
    \item $\mathbb{N}$ è limitato inferiormente ma non superiormente.
\end{itemize}
\end{esempio}

\begin{esercizio}
Quale tra i seguenti insiemi ha sia massimo sia minimo?
(1) $]2, 9[$ \qquad (2) $[0, 5[$ \qquad (3) $[0, 1[ \,\cup \{7\}$ \qquad (4) $]0, 2] \cup \{5\}$
\textbf{Risposta: 3.} $[0, 1[ \,\cup \{7\}$ ha minimo $0$ e massimo $7$. Gli altri: $]2, 9[$ non ha né minimo né massimo; $[0, 5[$ ha minimo $0$ ma non massimo; $]0, 2] \cup \{5\}$ ha massimo $5$ ma non minimo.
\end{esercizio}

\end{document}
